\section{Last Hope for the Dying}

If the party found \medicalHerb{} and returned it to \npcMedic{}, read the following:

\quoteBox{
	\npcMedic{}	rushes to his apparatus, brewing some medicine using freshly acquired herbs. Soon, he gives it to \npcSickGirl{}. She smiles calmly, pain gone from her face, but her eyes remain closed.
}

After characters tell \npcMedic{} about their last encounter with \enemyReanimatedAxehound{}, read the entry below.

If the party does not share this information after returning to camp, let \npcMedic{} ask about their wounds or other signs of the recent combat. Alternatively, find other opportunity to start a conversation about \enemyReanimatedAxehound{}s, that you can follow with:

\quoteBox{
	"I fear the girl has no much time... Nothing I do helps, and she has less and less life remaining in her with each passing hour... At this rate, she will not survive the night...
	\\ \\
	But those creatures you met give me a glimmer of hope... If those vines were able to strengthen dying axehounds enough to help them stand up and walk, even fight, perhaps some of this plant's properties could help us...?
	\\ \\
	But I would need a sample of those vines - fresh one, still animating some small creature... You can store it in this jar. Also, some more information would come in handy... Perhaps local populace would know more about them?"
}

The party is being sent to find the Red Vines samples. They can also learn more information about local flora in the Reshi village of the island.

If the party spent more time playing with \enemyAxehound{}s in the camp, \npcBabsk{} can propose sending a few of them with the characters on this mission. If not, a character may try to convince \npcBabsk{} to spare a few of them for the quest. 

\reasonBox{
	Depending on the approach the party would like to take to find the Red Vines samples, having a \enemyAxehound{} or two might help. It is not required though, and characters may venture without them.
}